\newpage
\section{Acta Número 02}
\label{sec:acta02}

\noindent \textbf{Fecha de la sesión:}\par
Lunes, 09 de marzo de 2026

\vspace*{0.5cm}
\noindent \textbf{Datos de la Sesión}
\vspace{-0.2cm}
\begin{itemize}[noitemsep]
    \item \textbf{Modalidad:} Virtual - Google Meet.
    \item \textbf{Hora de Inicio:} 18:00
    \item \textbf{Hora de Fin:} 20:00
\end{itemize}

\vspace*{0.5cm}
\noindent \textbf{Orden del Día}
\vspace{-0.2cm}
\begin{itemize}[noitemsep]
    \item Asistencia.
    \item Correciones al Documento de Invitación CPTIS-2302-2026
    \item Revisión de la Propuesta de Servicios: Análisis de los alcances y límites técnicos del sistema.
    \item Planificación y Plazo del Contrato: Aprobación del cronograma de Sprints y fechas límite.
    \item Análisis de la Propuesta Económica: Revisión de los costos de desarrollo y viabilidad financiera.
    \item Definición de Entregables y Documentación: Establecimiento de los artefactos y manuales a entregar al cliente.
    \item Firmas y Aprobación de los Presentes.
\end{itemize}

\vspace*{0.5cm}
\noindent \textbf{Socios Fundadores Presentes:}
\vspace{-0.2cm}
\begin{enumerate}[noitemsep]
    \item Pardo Pozo Juan Diego (Jefe de Proyecto).
    \item Arnez Jaimes Mauricio.
    \item Quispe Flores Camila Belen.
    \item Ariñez Bautista Jim Gabriel.
    \item Medina Rodríguez Mateo Jhoustin.
    \item Molina Maldonado Tomas Manuel.
    \item Zeballos Crespo Jonathan.
    \item Gutierrez Guzmán Angheline.
\end{enumerate}

\newpage
\subsection{Desarrollo del Acta}

\noindent \textbf{Control de Asistencia del Team}\par
Se procedió a la toma de asistencia a la hora acordada, corroborando la presencia en sesión de los 8 socios fundadores de la grupo-empresa Byte Busters S.R.L. Al contar con la totalidad de los miembros, se estableció el quórum necesario para iniciar la toma de decisiones y el trabajo de planificación.

\vspace*{0.5cm}
\noindent \textbf{Correciones al Documento de Invitación CPTIS-2302-2026}\par
Se revisó el documento de invitación a la convocatoria CPTIS-2302-2026, identificando y corrigiendo errores tipográficos, gramaticales y de formato. Se acordó la versión final del documento que será utilizado como referencia para la elaboración de las propuestas técnica y económica, asegurando su coherencia y profesionalismo. Además, se estableció un proceso de revisión interna para la generación de Anexos y Documentos de Soporte, garantizando la calidad y consistencia de toda la documentación relacionada con la convocatoria.\par

\vspace*{0.5cm}
\noindent \textbf{Revisión de la Propuesta de Servicios}\par
Se revisó exhaustivamente y se aprobó el documento correspondiente a la \textbf{Propuesta de Servicios}. Se establecieron claramente los objetivos generales y específicos, así como los alcances y límites técnicos del desarrollo de software requerido por la consultora TIS, delimitando la arquitectura a utilizar.

\vspace*{0.5cm}
\noindent \textbf{Planificación y Plazo del Contrato}\par
En base a los documentos de \textbf{Planificación} y \textbf{Plazo de Contrato}, el equipo definió el cronograma de trabajo estructurado en Sprints bajo la metodología SCRUM. Se acordaron las fechas de inicio y finalización de cada iteración, comprometiéndose los socios a cumplir estrictamente con los hitos de entrega estipulados en la convocatoria.

\vspace*{0.5cm}
\noindent \textbf{Análisis de la Propuesta Económica}\par
Se procedió a la lectura y validación de la \textbf{Propuesta Económica}. Todos los socios manifestaron su conformidad con la estimación de costos, incluyendo las tarifas por hora/hombre de los desarrolladores, los gastos operativos, de infraestructura en la nube y el presupuesto total detallado, asegurando la viabilidad y rentabilidad del proyecto.

\vspace*{0.5cm}
\noindent \textbf{Definición de Entregables y Documentación}\par
Finalmente, se revisó la sección de \textbf{Entregables y Documentación}. Se acordó la lista definitiva de los artefactos de software que serán proporcionados (código fuente en repositorios, scripts de base de datos) y se designaron responsabilidades para la consolidación de los manuales técnicos, de usuario y de instalación.

\vspace*{0.5cm}
\noindent \textbf{Aprobación Oficial}\par
Habiendo revisado la documentación de la convocatoria y tras aprobar unánimemente la Propuesta de Servicios, la Propuesta Económica y el cronograma de trabajo, se da por aprobada la sesión. Los miembros del equipo ratifican la viabilidad del proyecto y se comprometen a cumplir con los plazos y entregables estipulados.

\vspace{0.5cm}
\noindent \textbf{Firmas y Aprobación de los Presentes}
\begin{center}
    \noindent
    \setlength{\tabcolsep}{0pt}
    
    \begin{tabular}{ccc}
        \espacioFirma{assets/img/firmas/mauricio.jpg}{Mauricio Jaimes Arnez}{13751221} & 
        \espacioFirma{assets/img/firmas/camila.jpg}{Camila Belen Quispe Flores}{13097244} &
        \espacioFirma{assets/img/firmas/jim.jpg}{Jim Gabriel Ariñez Bautista}{9517642} \\[1.0cm]
        
        \espacioFirma{assets/img/firmas/mateo.jpg}{Mateo Medina Rodríguez}{9453809} &
        \espacioFirma{assets/img/firmas/tomas.jpg}{Tomas Molina Maldonado}{8051701} &
        \espacioFirma{assets/img/firmas/jonathan.jpg}{Jonathan Zeballos Crespo}{13067494} \\[1.0cm]
    \end{tabular}

    \vspace{0.2cm} 
    \begin{tabular}{cc}
        \espacioFirma{assets/img/firmas/diego.jpg}{Juan Diego Pardo Pozo}{12747783} & 
        \espacioFirma{assets/img/firmas/angheline.jpg}{Angheline Gutierrez Guzmán}{13751815}
    \end{tabular}
\end{center}

\vspace{0.5cm}
\noindent \textbf{Fecha de última revisión:} Lunes, 09 de marzo de 2026